\documentclass[11pt]{article}
\usepackage{amsmath,amssymb,amsthm,geometry}
\usepackage[hidelinks]{hyperref}
\geometry{margin=1in}
\newtheorem{proposition}{Proposition}
\newtheorem{remark}{Remark}
\newcommand{\R}{\mathbb R}
\newcommand{\dd}{\,\mathrm d}
\newcommand{\proj}{\operatorname{proj}}
\newcommand{\av}[1]{\left\langle #1\right\rangle}
\newcommand{\Cov}{\mathcal C}
\newcommand{\Cross}{\mathcal B}
\usepackage{mathrsfs}
\title{Actual coupled Navier--Stokes moments and quantitative source-chart orders\\
Retained third seed, old-wave interactions, and physical viscosity}
\author{}\date{}
\begin{document}\maketitle
This reader preserves the complete released-result propagation and actual
compact-curl moment calculation. It adds exact finite source-chart orders
for the retained third-growth seed and coefficient path. The source clock
$\nu=\sqrt{A_0}$ remains distinct from physical viscosity $\nu_{\rm NS}$;
all original parameters and signs are retained. The orders are finite-chart
identities with finite path constants, not uniform endpoint or infinite-cycle
estimates.
\tableofcontents\clearpage
\part{The actual released result, propagation, and radial moments}
\section{Exact operator bridge to the released Navier--Stokes profiles}
\label{nsbridge:nsb:operator}

This bounded comparison uses the locally saved released Navier--Stokes source
(the source pages are 7--15 and 24--27 for the profile, stress and pulse
construction). Its stated physical viscosity is denoted by $\nu_{\rm NS}>0$.
The Boussinesq lane variables and physical viscosity remain $d$, $\mu$ and
$\nu=\sqrt{A_0}$; none is identified with $\nu_{\rm NS}$.

Use right-handed cylindrical coordinates $(r,\vartheta,z)$ and define
\[
 y_1=z,\qquad y_2=\frac{r^2}{2}>0,\qquad
 v=(u^z,ru^r)=J\nabla_y\Psi,
 \quad J=\begin{pmatrix}0&-1\\1&0\end{pmatrix}.
\]
The meridional divergence identity is
$\partial_{y_1}v_1+\partial_{y_2}v_2=0$. Define the circulation and
meridional vorticity variables
\[
 \Gamma=ru^\vartheta,
 \qquad \xi=\frac{\omega^\vartheta}{r},
 \qquad h=\frac{\Gamma}{2y_2}=\frac{u^\vartheta}{r},
 \qquad Q=\partial_{y_1}^2+2y_2\partial_{y_2}^2,
 \qquad M=Q+4\partial_{y_2}.
\]
The azimuthal momentum and meridional vorticity rows of the cylindrical
equations (with arbitrary smooth force components) become, on $y_2>0$,
\begin{align}
 D_v\Gamma&=\nu_{\rm NS}Q\Gamma+r f^\vartheta,
 \label{nsbridge:nsb:Gamma}\\
 D_v\xi&=\nu_{\rm NS}M\xi+\partial_{y_1}(h^2)
       +\operatorname{curl}_y(f^z,f^r/r),
 \label{nsbridge:nsb:xi}
\end{align}
where $D_v=\partial_t+v\cdot\nabla_y$ and
$\operatorname{curl}_y(a,b)=\partial_{y_1}b-\partial_{y_2}a$.

Indeed, $\partial_r=r\partial_{y_2}$ and
$\partial_r^2+r^{-1}\partial_r=2y_2\partial_{y_2}^2+2\partial_{y_2}$.
For $\Gamma$, the azimuthal viscous operator is
$\partial_r^2-r^{-1}\partial_r+\partial_z^2=Q$.
For $r\xi=\omega^\vartheta$, the cylindrical operator
$(\partial_r^2+r^{-1}\partial_r-r^{-2}+\partial_z^2)(r\xi)$,
divided by $r$, is $M\xi$; explicitly the radial terms are
$2y_2\xi_{y_2y_2}+4\xi_{y_2}$. The centrifugal term is
$r^{-2}\partial_z((u^\vartheta)^2)=\partial_{y_1}(\Gamma^2/(4y_2^2))
=\partial_{y_1}(h^2)$. These calculations prove every coefficient in
(\ref{nsbridge:nsb:Gamma})--(\ref{nsbridge:nsb:xi}), including the $r$ force conversion.

The velocity--vorticity relation is also retained. Since
$u^z=-\Psi_{y_2}$ and $u^r=\Psi_{y_1}/r$,
\[
 \xi=\frac{\partial_z u^r-\partial_r u^z}{r}
     =\frac{\Psi_{y_1y_1}}{2y_2}+\Psi_{y_2y_2}
     =\frac{Q\Psi}{2y_2}.
\]
Thus the two transport rows do not by themselves replace the elliptic
velocity reconstruction. Conversely, on a simply connected meridional
domain, a smooth $\Psi,\Gamma$ satisfying this relation and both transport
rows recovers $u$. Define its momentum residual without pressure as
$a=\partial_tu+(u\cdot\nabla)u-\nu_{\rm NS}\Delta u-f$.
The circulation row is $a^\vartheta=0$, while the vorticity row is
$\partial_z a^r-\partial_r a^z=0$. The one-form
$a^r\,\mathrm dr+a^z\,\mathrm dz$ is therefore closed, and
$p=-\int(a^r\,\mathrm dr+a^z\,\mathrm dz)$ along a path from a fixed
base point is path independent. Its derivatives are $p_r=-a^r$,
$p_z=-a^z$. This proves the local momentum reconstruction, with pressure
unique up to a function of time; boundary data are not removed.

The change from circulation to swirl rate is invertible on the exact retained
domain: $\Gamma=2y_2h$. Product differentiation gives
\[
 Q(2y_2h)=2y_2Mh,\qquad
 D_v(2y_2h)=2y_2D_vh+2v_2h.
\]
Therefore (\ref{nsbridge:nsb:Gamma}) is equivalent to
\begin{equation}
 D_vh=\nu_{\rm NS}Mh-\frac{v_2}{y_2}h+\frac{f^\vartheta}{r}.
 \label{nsbridge:nsb:h}
\end{equation}
The source profile uses precisely this relation: in its similarity variables,
$E$ is the azimuthal velocity profile and
$H_{\rm src}=\sqrt{2X}E$ is the scaled circulation profile (source equations
(4.3), (4.8)). Since $r=\sqrt{2qX}$ and
$u^\vartheta=q^{-A}E$, its physical circulation is
$\Gamma=q^{1/2-A}H_{\rm src}=q^{-h_{\rm src}}H_{\rm src}$,
where the source exponents satisfy $A=1/2+h_{\rm src}$. Here $q>0$ is the
geometric similarity coordinate and $h_{\rm src}>0$ is an exponent; neither
is a signed amplitude scale. This is the exact scale dictionary; it does not
identify $H_{\rm src}$ with the Boussinesq scalar amplitude.

The sign of the blowout amplitude must be kept separate from that positive
geometric coordinate. The axisymmetric equations have the exact involution
\[
 (u^r,u^\vartheta,u^z,p;f^r,f^\vartheta,f^z)
 \longmapsto
 (u^r,-u^\vartheta,u^z,p;f^r,-f^\vartheta,f^z).
 \label{nsbridge:nsb:sign}
\]
It sends $(\Gamma,h)$ to $(-\Gamma,-h)$, leaves $\xi$ and the meridional
flow unchanged, and preserves every displayed equation: the only quadratic
swirl term is $\partial_{y_1}(h^2)$, while the azimuthal equation is linear in
$(\Gamma,f^\vartheta)$. Consequently a source-selected positive profile
$E>0$ (source pp. 7--9) does not justify assuming that every blowout branch is
positive; its reflected branch has $E<0$, $\Gamma<0$ and the same blowup
magnitude. The exponent $q^{-h_{\rm src}}$ controls the magnitude in either
branch. The source's positive covariance coefficients and viscosity are
separate construction choices, not a sign claim about the velocity amplitude.

For comparison with the lane's controlled scalar equation, set
$B=h^2$. The globally smooth square equation, including its geometric and
viscous terms, is
\begin{equation}
 D_vB=\nu_{\rm NS}MB-2\nu_{\rm NS}
 |\nabla h|_Q^2-\frac{2v_2}{y_2}B+\frac{2h}{r}f^\vartheta,
 \qquad |\nabla h|_Q^2=h_{y_1}^2+2y_2h_{y_2}^2.
 \label{nsbridge:nsb:B}
\end{equation}
This follows by multiplying (\ref{nsbridge:nsb:h}) by $2h$ and using
$M(h^2)=2hMh+2|\nabla h|_Q^2$. At zeros of $h$ equation
(\ref{nsbridge:nsb:B}) remains smooth; dividing by $B$ is unnecessary and would lose
that domain. On a component where $B>0$, the equivalent expression is
$D_vB=\nu_{\rm NS}MB-\nu_{\rm NS}|\nabla B|_Q^2/(2B)
-2v_2B/y_2+2hf^\vartheta/r$.

The exact relation to the Boussinesq lane is consequently an operator
correspondence, not a scalar conjugacy. The lane's retained-viscosity scalar
row has $D_v\Theta=d\Delta\Theta+f_\Theta$ in Cartesian material
coordinates. A direct identification $\Theta=B$ would require, on the same
physical domain, both $M=\Delta$ after the nonlinear coordinate pullback and
$-2v_2B/y_2-2\nu_{\rm NS}|\nabla h|_Q^2+2hf^\vartheta/r$ to be absorbed into
one prescribed force. The first condition fails as an operator identity:
$M=\partial_{y_1}^2+2y_2\partial_{y_2}^2+4\partial_{y_2}$ has variable metric
and drift, whereas the lane's $\Delta$ is the Euclidean Cartesian Laplacian.
The geometric, gradient and force terms must therefore be retained or
calculated in the particular source state. Nonconstant profiles alone
would not prove that such terms never cancel.
The exact bridge is (\ref{nsbridge:nsb:Gamma})--(\ref{nsbridge:nsb:B}) with the
full residual terms retained. It neither claims that the lane's controls are
the source pulses nor changes the source's pressure reconstruction.

For the meridional velocity, the source's pulse construction supplies the
following further exact local dictionary (source pp. 11--15 and 74--77). If
$w$ is a divergence-free pulse and $\pi$ its pressure increment, then
\[
 R(u_B+w,p_B+\pi)=R(u_B,p_B)+L_{u_B}(w,\pi)
 +\nabla\!\cdot(w\otimes w),
\]
\[
 L_{u_B}(w,\pi)=\partial_tw+(u_B\cdot\nabla)w
 +(w\cdot\nabla)u_B-\nu_{\rm NS}\Delta w+\nabla\pi,
 \qquad \nabla\cdot w=0.
\]
For the moving local phase with wavevector $n$ and transverse amplitude
$t_m$, the source equation is
\[
 t_m'+\mathcal Kt_m+m^2d_N t_m+i k m n\,\pi_m=-f_m,
 \qquad n\cdot t_m=0,
 \qquad d_N=\varepsilon k^2|n|^2,
\]
with the source's moving shear matrix
$\mathcal K=\begin{pmatrix}0&-2F&0\\2F+RF_R&0&0\\G_R&0&0\end{pmatrix}$.
The transverse projection uses
\[
 A_\Phi=-\mathcal K+
 \frac{n(n^T\mathcal K-n'^T)}{|n|^2}.
\]
These are exact source pulse equations, including viscosity and pressure
projection. Their quadratic covariance is chosen to cancel the annular
stress; the source explicitly retains curl, cutoff, and correction interactions.
The Boussinesq lane's two-dimensional pair and its full profile residuals have
an exact algebraic form, but no identification with this three-dimensional
pulse equation follows without an additional embedding that supplies the
source's cylindrical geometry, moving plane, covariance cone and correction
cycle. That embedding is not asserted here.

The current bridge therefore proves a precise morphism of variables, operators,
forces and pulse residuals and records the exact obstruction to a direct scalar
identification. It establishes the next programme datum while preserving the
unproved status of any singular limit or third return.


\sloppy

\section{Scope and source equations}
This note records one local, exact comparison.  It does not identify the
coupled Boussinesq stages with the released Navier--Stokes construction and
makes no claim about a global Navier--Stokes solution.  The source is the
released PDF \texttt{openai\_navier\_stokes.pdf}, SHA-256
\texttt{8c8a94ad9ac824c8b605b9827cadf7beaca48bd10b380de3cfc872a2c37afa81}.
The source statements used here are Section 7, equations (7.2)--(7.8),
(7.12)--(7.22), and Section 9, equations (9.1)--(9.2), on PDF pages 74--75,
77--81, and 101--103. In the source notation, on one lifted pulse rectangle,
\begin{equation}
 \Phi=p\theta+p_zZ/\varepsilon+x_0R-v(pF+p_zG),\qquad
 n=\nabla_*\Phi,
\end{equation}
with
\begin{equation}
 n=(x_0-v(pF_R+p_zG_R),\ p/R,\ p_z-\varepsilon v(pF_Z+p_zG_Z)).
 \label{nsbridge:eq:nsource}
\end{equation}
For a nonzero integer harmonic $m$, the exact principal pulse equations are
\begin{equation}
 n\cdot t_m=0,\qquad
 t_m'+\mathsf Kt_m+m^2d_Nt_m+ikm n\pi_m=-f_m,
 \qquad d_N=\varepsilon k^2|n|^2,
 \label{nsbridge:eq:sourcepulse}
\end{equation}
where a prime is the pulse-coordinate derivative and
\begin{equation}
 \mathsf K=\begin{pmatrix}0&-2F&0\\2F+RF_R&0&0\\G_R&0&0\end{pmatrix},\qquad
 \mathsf A_\Phi=-\mathsf K+\frac{n(n^T\mathsf K-(n')^T)}{|n|^2}.
 \label{nsbridge:eq:sourceoperator}
\end{equation}
The source pressure coefficient is exactly
\begin{equation}
 \pi_m=\frac{i}{km|n|^2}\bigl(n\cdot\mathsf Kt_m-(n')\cdot t_m+n\cdot f_m\bigr).
 \label{nsbridge:eq:sourcepressure}
\end{equation}
For $m=1$, a real homogeneous pulse $t_1$ has the source energy identity
(7.22):
\begin{equation}
 \frac12(|t_1|^2)'=-g\cdot\bigl(t_{1,r}(t_{1,\theta},t_{1,z})\bigr)
 -\varepsilon k^2|n|^2|t_1|^2,
 \label{nsbridge:eq:sourceenergy}
\end{equation}
where $g=(RF_R,G_R)$ in the tangential coordinates.  Thus the source pulse
retains both shear transfer and viscous dissipation.

\section{The coupled physical pair}
For one activated layer of the coupled calculation, put
\begin{equation}
 a_c=-\frac{J\zeta\cdot G_{<q}}{Kr^2},\qquad b_c=K\zeta_1,\qquad
 \lambda_c^\Theta=\kappa K^2r^2,\qquad \lambda_c^\Omega=\varepsilon_c K^2r^2.
\end{equation}
Its exact retained-viscosity pair, including optional scalar and vorticity
controls, is
\begin{equation}
 \dot y=C_cy+c_c,\qquad y=\binom{\Theta}{\Omega},\qquad
 C_c=\begin{pmatrix}-\lambda_c^\Theta&a_c\\b_c&-\lambda_c^\Omega\end{pmatrix},\quad
 c_c=\binom{c_\theta}{c_\omega}.
 \label{nsbridge:eq:coupledpair}
\end{equation}
Here $G_{<q}$, $\zeta$, and $r$ evolve with the inherited background and phase
ODEs; in particular no fixed-background substitution is made.  Let $v$ be a
source pulse coordinate on an interval and let $q_v=\dd t/\dd v$ describe
its physical-time parametrization. The forward source parametrization has
$q_v>0$; the algebra below also retains a negative orientation without
changing the amplitude sign. Define $y'(v)=q_v(C_cy+c_c)$.
Here $K,r$ are the coupled frequency and phase length, not the source
carrier $k$ and cylindrical radius $R$. The original $A_0,\lambda_0,\mu$
and time scale $\nu=\sqrt{A_0}$ are unchanged. The source chart parameter
$\varepsilon$ is not identified with either physical diffusivity.

\section{Exact local residual map}
Let $B(v)\in\R^{3\times2}$ be a $C^1$ tangent frame with full column rank and
$n(v)^TB(v)=0$, so its range is exactly $n(v)^\perp$.  Let $B_\ell(v)$ be any
left inverse, $B_\ell B=I_2$; then $B B_\ell$ acts as the identity on
$n^\perp$.  Let $L(v)\in\R^{2\times2}$ be any $C^1$ coefficient map and embed
the coupled pair into a transverse source amplitude by
\begin{equation}
 t(v)=B(v)L(v)y(v).
 \label{nsbridge:eq:embedding}
\end{equation}
For an embedding with a non-tangent frame the transversality condition would be
\begin{equation}
 n(v)^TB(v)L(v)y(v)=0.
 \label{nsbridge:eq:transversecondition}
\end{equation}
For the tangent frame used here, the condition holds for every $L,y$.  Set
\begin{equation}
 H(v)=B_\ell(v)\bigl(\mathsf A_\Phi(v)B(v)-B'(v)\bigr).
 \label{nsbridge:eq:Hdef}
\end{equation}
Then the exact projected source residual of the embedded coupled pair is
\begin{align}
 \mathcal E_{L,c}:=B_\ell\proj_{n^\perp}
 \left[t'-\mathsf A_\Phi t+m^2d_Nt+f_m\right]
 &=\left[L'+q_vLC_c-HL+m^2d_NL\right]y+q_vLc_c \\
 &\quad+B_\ell\proj_{n^\perp}f_m.
 \label{nsbridge:eq:residualmap}
\end{align}
In particular, with $L=I$ and no added source, the defect is the explicit
matrix expression
\begin{equation}
 \mathcal E_{I,0}=\bigl(q_vC_c-H+m^2d_NI_2\bigr)y.
 \label{nsbridge:eq:identitydefect}
\end{equation}
This gives a local residual relation with every actual coefficient retained.

\begin{proof}
Differentiate \eqref{nsbridge:eq:embedding}:
$t'=B'L y+B L'y+B L y'$.  Since $y'=q_v(C_cy+c_c)$,
differentiate $n^TB=0$ to obtain $n^TB'=-(n')^TB$. Also
$n^T\mathsf A_\Phi B=-(n')^TB$, directly from
\eqref{nsbridge:eq:sourceoperator}. Thus $\mathsf A_\Phi B-B'$ has transverse
range and equals $BH$. Applying $B_\ell$ to the transverse terms and
$B_\ell\operatorname{proj}_{n^\perp}$ to the force proves
\eqref{nsbridge:eq:residualmap}. More explicitly,
\begin{equation}
 n^T(t'-\mathsf A_\Phi t+m^2d_Nt)=0,\qquad
 n^T(t'-\mathsf A_\Phi t+m^2d_Nt+f_m)=n^Tf_m.
\end{equation}
The unprojected force need not be transverse. The pressure below removes
precisely this remaining normal component.
\end{proof}

\section{Pressure reconstruction and full three-dimensional residual}
For any embedded transverse $t$ and arbitrary $f_m$, define
\begin{equation}
 \pi_{L,c}=\frac{i}{km|n|^2}
 \left(n\cdot\mathsf Kt-(n')\cdot t+n\cdot f_m\right).
 \label{nsbridge:eq:pressuremap}
\end{equation}
Then the full source residual satisfies the exact decomposition
\begin{align}
 \mathcal R_{L,c}
 &: =t'+\mathsf Kt+m^2d_Nt+ikm n\pi_{L,c}+f_m \\
 &=\proj_{n^\perp}\left(t'+\mathsf Kt+m^2d_Nt+f_m\right)
 =B\,\mathcal E_{L,c},
 \label{nsbridge:eq:fullresidual}
\end{align}
where the last equality uses the definition of $H$ and the transverse condition.
Thus the pressure removes precisely the normal component; it does not remove
$\mathcal E_{L,c}$, the moving-frame, shear, diffusion, or coupled-background
mismatch.  If one elects to force the source pulse to obey the released equation,
then the unique additional transverse forcing is
\begin{equation}
 f_{m,\mathrm{corr}}=-B\,\mathcal E_{L,c},
\end{equation}
added to the old $f_m$. The old $\mathcal E_{L,c}$ is evaluated before
this addition, so the new projected residual is zero.
Since $n^Tf_{m,\mathrm{corr}}=0$, pressure is unchanged.
This is an exact local
forced relation, not an identification of the physical constructions.

\section{Diffusion and energy comparison}
The source diffusion multiplier is $m^2d_N=m^2\varepsilon k^2|n|^2$ in
\eqref{nsbridge:eq:sourcepulse}.  The coupled pair has the two retained physical losses
$\lambda_c^\Theta$ and $\lambda_c^\Omega$, which are generally unequal to
$m^2d_N$.  Their signed discrepancy in \eqref{nsbridge:eq:identitydefect} is explicitly
\begin{equation}
 \bigl(m^2d_N-q_v\lambda_c^\Theta\bigr)\Theta,
 \qquad
 \bigl(m^2d_N-q_v\lambda_c^\Omega\bigr)\Omega
\end{equation}
when $H=0$ and $C_c$ is expanded entrywise; with general $H$ the exact
diagonal entries are those of $q_vC_c-H+m^2d_NI_2$.  No damping is dropped.
If one compares nonnegative loss magnitudes instead, the corresponding cost
gaps are $|m^2d_N-q_v\lambda_c^\Theta|$ and
$|m^2d_N-q_v\lambda_c^\Omega|$.
The required finite-interval coordinate map can be constructed explicitly.
Let $I=[v_0,v_1]$, and let $X_c,X_s$ be the identity-initial-value
fundamental matrices for $D_c=q_vC_c$ and $D_s=H-m^2d_NI_2$,
respectively. For either continuous matrix $D$, define
\[
 X_D(v)=I_2+\sum_{j\ge1}
 \int_{v_0\le w_j\le\cdots\le w_1\le v}
 D(w_1)\cdots D(w_j)\,\dd w_j\cdots\dd w_1 .
\]
With $M_D=\sup_I\|D\|$ and $T_I=v_1-v_0$, the $j$th integral has norm
at most $(M_DT_I)^j/j!$. Uniform convergence proves the integral equation
$X_D=I_2+\int_{v_0}^vDX_D$ and hence $X_D'=DX_D$.
The same iteration for $Z'=-ZD$, $Z(v_0)=I_2$ gives $ZX_D=I_2$
by differentiation, proving invertibility on $I$. For a specified
invertible $L_0$, set
\[
 L=X_sL_0X_c^{-1},\qquad L^{-1}=X_cL_0^{-1}X_s^{-1}.
\]
This has norm at most $\|L_0\|e^{(M_c+M_s)T_I}$; its inverse has the
same bound with $L_0^{-1}$. Differentiation proves
\begin{equation}
 L'=HL-q_vLC_c-m^2d_NL,
 \label{nsbridge:eq:intertwiner}
\end{equation}
so the homogeneous embedded pair has zero projected residual. This proves
an actual finite-interval conjugacy along the retained coefficient path.
It does not make the source phase,
pressure, cutoff, or nonlinear covariance equal to the coupled Boussinesq fields.
For general controls $c_c$, the forced term is $q_vLc_c$ in
\eqref{nsbridge:eq:residualmap}.

For an arbitrary embedded field (possibly complex), the exact energy balance
retains its full defect:
\begin{equation}
 \begin{split}
 \frac12(|t|^2)'={}&-\operatorname{Re}(\overline t^{\,T}\mathsf Kt)
 -m^2d_N|t|^2-\operatorname{Re}(\overline t^{\,T}f_m)\\
 &+\operatorname{Re}(\overline t^{\,T}B\mathcal E_{L,c}).
 \end{split}
\end{equation}
This follows by taking the real Hermitian product of
\eqref{nsbridge:eq:fullresidual} with $t$. Its pressure term vanishes because $n$ is
real and $n^Tt=0$. The $2F$ part of $\mathsf K$ is skew-symmetric, hence
its real work is zero; the shear work equals
$\operatorname{Re}(\overline{t_r}(g_\theta t_\theta+g_z t_z))$.
Only for a zero residual, zero force, $m=1$, and real $t$ does this reduce
to \eqref{nsbridge:eq:sourceenergy}. The real coupled pair has
\begin{equation}
 \frac12\frac{\dd}{\dd t}|y|^2
 =-\lambda_c^\Theta\Theta^2-\lambda_c^\Omega\Omega^2
 +(a_c+b_c)\Theta\Omega+c_\theta\Theta+c_\omega\Omega,
 \label{nsbridge:eq:coupledenergy}
\end{equation}
which is a two-component physical energy identity.  Equations
\eqref{nsbridge:eq:residualmap} and \eqref{nsbridge:eq:coupledenergy} exhibit exactly where the
source's three-dimensional shear flux and the coupled pair's two-dimensional
exchange differ.

\section{Proven scope}
The source uses a cylindrical, localized, divergence-free curl construction
(Section 7.4 and Section 9, PDF pages 85--87 and 102--103); the displayed
principal pulse equation is only one component of that construction.  The
comparison above proves an exact local map and its full pressure-projected
residual for every finite interval on which the denominators and frame left
inverse exist.  It preserves source viscosity, moving phase normal, pressure,
rotation/shear matrix, and all coupled physical diffusion coefficients.  It does
not prove that the coupled third-growth or steering stages produce the source's
pulse phase $\Phi$, covariance, cutoff tails, or smooth terminal force.  Those
remain separate bridge problems.

The sign conventions in this comparison are also explicit.  The quantities
$q_v>0$, $\varepsilon>0$, and $d_N=\varepsilon k^2|n|^2>0$ record the chosen
orientation of physical time and the positive viscous coefficient; they are
not assumptions on the sign of a pulse or blowout amplitude.  The equations
\eqref{nsbridge:eq:sourcepulse}--\eqref{nsbridge:eq:sourcepressure} are linear in $(t_m,f_m)$,
so the simultaneous replacement $(t_m,f_m,\pi_m)\mapsto
(-t_m,-f_m,-\pi_m)$ is an exact reflected branch.  Likewise, the local map
\eqref{nsbridge:eq:embedding} is sign-equivariant under $y\mapsto-y$ and
$c_c\mapsto-c_c$, with the coefficients fixed at the actual path.
This is a statement about the retained linear equations, not a sign symmetry
of the entire coupled nonlinear system.
The source covariance is quadratic, so reversing a wave leaves that
covariance unchanged. Reversing the full Navier--Stokes velocity need not
reverse its force, since $(u\cdot\nabla)u$ is even in $u$.
The full-flow map that reverses axisymmetric swirl is a spatial reflection;
its complete formula is proved in the accompanying covariance bridge.
No sign of a blowout or endpoint is inferred from positive geometric scales.


\section{Source objects and the finite calculation}
This calculation uses the preserved 165-page supplied manuscript
\emph{Finite Time Blowup for Navier--Stokes}, source equations
(7.23)--(7.42), (9.12)--(9.14). Its PDF SHA-256 is
\begin{center}\small
\texttt{8c8a94ad9ac824c8b605b9827cadf7beaca48bd10b380de3cfc872a2c37afa81}.
\end{center}
The source pulse fields, auxiliary rectangles, phases, and leading
covariance columns are the given objects; their global construction is
source input. Every finite map asserted below is proved here. The existing
coupled-stage proofs remain in the complete 102-page reader. The companion
pulse note proves the exact frame residual and a constructed finite-interval
intertwiner along its evolving coefficients. The source uses positive
$q(z,t)$, band scales $Q_\beta=2^{-\ell_\beta}$,
$\varepsilon_\beta=Q_\beta^h$, $A=1/2+h$; none is set to one here or
identified with the original coupled $A_0,\lambda_0,\mu,\nu=\sqrt{A_0}$
or its physical diffusivities.

All operations in this note are on a compact set inside $q>0$ and the
source active annulus. There are finitely many relevant bands and
locally finite source labels. Supports of newly prescribed inputs are
strictly inside the annulus, so no lower bound on a leading amplitude
at a vanishing shell edge is assumed. Endpoint and shell-edge estimates
are not inferred from compact-interior estimates.

\section{Physical curl of the actual retained amplitude}
Let $y=(\Theta,\Omega)^T$ be the actual smooth retained coupled solution,
$y'=q_v(C_cy+c_c)$, with the coefficients of the companion pulse note.
Use the smooth actual coefficients on the chosen slow/auxiliary patch;
all harmonic coefficients and cutoffs are independent of $\theta$ in
cylindrical components, and each carrier has nonzero integer angular
frequency. For its source tangent frame $B$, smooth coefficient map $L$,
and a fixed smooth real compact cutoff $\chi$ of this type, put
\[
 t_m=\chi BLy,\qquad
 \mathcal E_\chi=
 \chi\{[L'+q_vLC_c-HL+m^2d_NL]y+q_vLc_c\}+\chi'Ly,
 \quad H=B_\ell(\mathsf A_\Phi B-B').
\]
Here prime differentiates along the pulse coordinate; the other slow
and auxiliary derivatives remain in the spatial formulas below.
Smooth parameter dependence of the companion's constructed $L$ follows
directly from its integral equations on a fixed interval: for a nonzero
parameter multi-index $I$,
\[
 \partial^I X(v)=
 \int_{v_0}^v\left[D\,\partial^I X+
 \sum_{0<J\le I}\binom{I}{J}(\partial^JD)
                                 \partial^{I-J}X\right]\,\dd w .
\]
The initial parameter derivatives vanish for identity initial data.
Inductively the known forcing is smooth and bounded on the compact patch;
the same convergent ordered-integral iteration solves this equation.
This proves every finite mixed derivative of $X,X^{-1},L$.
Since $n^TB=0$, $n^Tt_m=0$. Differentiation proves
$t_m'-\mathsf A_\Phi t_m+m^2d_Nt_m=B\mathcal E_\chi$.
Thus even the pulse-cutoff term is explicit.

A physical version of the exact curl construction makes the intermediate
map to momentum residual unambiguous. For each chosen harmonic let
$\phi(x,t)=k_\beta m\Phi_\beta(x,t,Y(x,t))$ be its real evaluated phase,
let $\xi=\nabla_x\phi\ne0$, and let $a(x,t)=Q_\beta^{-A}t_m$ be the
physical transverse amplitude. The source coordinate change is
$r=\sqrt{Q_\beta}R$, $z=Q_\beta^D Z$, $1-t=Q_\beta T$,
with $D=1/2-h$ and normalized axial derivative
$D_z=\varepsilon_\beta\partial_Z$. It gives
$\xi=(k_\beta m/\sqrt{Q_\beta})n_\Phi$; all auxiliary chain derivatives
are included in $n_\Phi$. Hence $\xi\cdot a=0$. Define in Cartesian
components
\begin{equation}
 C=\frac{i\,\xi\times a}{|\xi|^2},\quad
 \mathcal A=C e^{i\phi},\quad
 \nabla\times\mathcal A=(a+r_a)e^{i\phi},\qquad
 r_a=\nabla\times C. \label{nsbridge:cb:curl}
\end{equation}
The vector triple-product identity gives
$i\xi\times C=a$, proving the last equality. It also shows that the
physical potential has exactly the source factor $Q_\beta^{1/2-A}$
relative to its chart coefficient
$i\,n_\Phi\times t_m/(k_\beta m|n_\Phi|^2)$.
Cartesian differentiation of $C$ includes derivatives of the cylindrical
basis; the same terms are the $R^{-1}$ terms of source (7.37)--(7.38).
Take real parts and sum the finite retained harmonics. This constructs
a smooth compact velocity $w_y$ with $\nabla\cdot w_y=0$.

Retain a finite smooth source state $(u_*,p_*)$, including its background,
actual leading waves and their curl corrections, and the real
physical pressure increments obtained from the companion pulse pressure
formula with $f_m=0$:
\[
 \pi_{\chi,m}=\frac{i}{k_\beta m|n|^2}
                  \bigl(n^T\mathsf Kt_m-(n')^Tt_m\bigr),\qquad
 \pi_y=\sum_{\beta,m}\operatorname{Re}
             (Q_\beta^{-2A}\pi_{\chi,m}e^{i\phi_{\beta,m}}).
\]
These are smooth and compact wherever $t_m$ is. This definition fixes the
pressure map without assuming the pulse residual vanishes.
Define the actual candidate residual
\begin{align}
 R_y={}&R_{\nu_{\rm NS}}(u_*+w_y,p_*+\pi_y), \nonumber\\
 R_{\nu_{\rm NS}}(u,p)
   :={}&\partial_tu+(u\cdot\nabla)u-\nu_{\rm NS}\Delta u+\nabla p .
 \label{nsbridge:cb:Ry}
\end{align}
This is a definition by explicit differentiated fields, not an assumed
source term. Its exact expansion is
\[
 R_y=R_{\nu_{\rm NS}}(u_*,p_*)
 +\partial_tw_y+(u_*\cdot\nabla)w_y+(w_y\cdot\nabla)u_*
 -\nu_{\rm NS}\Delta w_y+\nabla\pi_y
 +(w_y\cdot\nabla)w_y .
\]
Every cutoff, parent, phase, and curl derivative occurs in this formula.
Physical $\nu_{\rm NS}$ is retained and is distinct from every source
time-scaling convention. Define
$\Delta R_y=R_y-R_{\nu_{\rm NS}}(u_*,p_*)$.
We correct the new compact increment; this difference has compact support
since every added term does. The baseline residual remains explicitly.

The averaging needed next is the source extended-field operation, not a
spatial average of the evaluated field. Evaluate derivatives on the
extended coefficients first by the complete chain rule. In a fixed
chart set $\mathcal D_i=\partial_{x_i}
+\sum_\alpha(\partial_{x_i}Y_\alpha)\partial_{Y_\alpha}$ and similarly
$\mathcal D_t$. These are the pullbacks of commuting physical
derivatives, so $[\mathcal D_i,\mathcal D_j]=0$ on their lifted domain.
For smooth periodic coefficients their auxiliary derivative terms have
zero Haar average. In cylindrical components angular differentiation
also differentiates the basis. This defines an extended residual
$\widetilde R_y$ whose evaluation is exactly \eqref{nsbridge:cb:Ry}.
Average the tangential components of the increment at fixed slow variables:
\begin{equation}
 F_2=\av{\Delta\widetilde R_{y,\theta}},\qquad
 F_1=\av{\Delta\widetilde R_{y,z}},\qquad
 \av{g}=\int_{\mathbb T^2}\frac1{2\pi}
                    \int_0^{2\pi}g\,\dd\theta\,\dd Y . \label{nsbridge:cb:F}
\end{equation}
These $F_e$ are smooth and compact
in the chosen physical annulus. Equations \eqref{nsbridge:cb:curl}--\eqref{nsbridge:cb:F}
are the explicit intermediate map from the retained coupled pair to a
signed momentum residual. They do not identify it with the original
two-dimensional Boussinesq force.

\section{Radial moments, support, and the original chart scales}
For each fixed $(z,t)$, choose $0<r_-<r_+$ containing the compact
supports of both $F_e$ and the chosen bumps in the active annulus.
The source bump is
\[
 b_e(r,z,t)=q^{-(e+1)/2}\widehat b_e(r/\sqrt q),\qquad
 \int_0^\infty x^e\widehat b_e(x)\,\dd x=1,\qquad e=1,2,
\]
with interior support. The change $r=\sqrt q\,x$ proves
$\int r^eb_e\,\dd r=1$ exactly. Define
\begin{equation}
 M_e=\int_0^\infty r^eF_e(r)\,\dd r,\qquad
 \sigma_e(r)=-r^{-e}\int_0^r s^e(F_e(s)-b_e(s)M_e)\,\dd s .
 \label{nsbridge:cb:primitive}
\end{equation}
Then
\begin{equation}
 (\partial_r+e/r)\sigma_e=-F_e+b_eM_e,\qquad
 \int_0^\infty r^e(F_e-b_eM_e)\,\dd r=0. \label{nsbridge:cb:radial}
\end{equation}
Differentiate $r^e\sigma_e$ to prove the first identity; integrate and
use the bump normalization for the second. The integral is zero below
$r_-$ and above $r_+$, so $\sigma_e$ has compact smooth support in the
same enclosing interval. A homogeneous solution is $cr^{-e}$, and a
compact homogeneous solution has $c=0$; thus the primitive is unique
among compact solutions of its displayed equation. Moreover a compact
solution of $(\partial_r+e/r)\sigma=-F_e$ exists if and only if $M_e=0$:
necessity follows by integrating the weighted derivative, and
sufficiency follows from \eqref{nsbridge:cb:primitive}. The two moments are
therefore the exact surviving finite-dimensional terms.

For $C_e=\int_{r_-}^{r_+}s^e\,\dd s$,
$|M_e|\le C_e\|F_e\|_\infty$ and
\[
 \|\sigma_e\|_\infty\le r_-^{-e}C_e
        (\|F_e\|_\infty+\|b_e\|_\infty|M_e|).
\]
If $G_e=F_e-b_eM_e$, successive radial derivatives obey the exact recurrence
\[
 \partial_r^{j+1}\sigma_e=-\partial_r^jG_e
 -e\sum_{i=0}^j \binom{j}{i}(-1)^ii!r^{-i-1}
                              \partial_r^{j-i}\sigma_e .
\]
This proves every finite radial derivative bound inductively. Parameter
derivatives follow by differentiating the compact integral; all derivatives
of $q,b_e$ and the moving physical data are retained. On the present compact
set $q$ is bounded away from zero. No terminal uniform estimate is implied.

The source (9.12)--(9.14) uses physical
$F_e=Q^{-2A-1/2}E_e^*$, where
$E_2^*=\av{E_\theta}_Y$ and $E_1^*=\av{E_z}_Y$.
For our fields define $E_e^*=Q^{2A+1/2}F_e$ in that same convention.
With $r=\sqrt Q R$ and $Q$ fixed in a chart, the exact conversions are
\begin{gather}
 M_e=Q^{e/2-2A}M_e^*,\quad M_e^*=\int R^eE_e^*\,\dd R,\quad
 b_e^*(R)=Q^{(e+1)/2}b_e(\sqrt Q R),\nonumber\\
 \Sigma_e(R)=Q^{2A}\sigma_e(\sqrt Q R),\qquad
 (\partial_R+e/R)\Sigma_e=-E_e^*+b_e^*M_e^*. \label{nsbridge:cb:chart}
\end{gather}
Substitute $r=\sqrt Q R$ including its measure in the moment, and
differentiate $\Sigma_e$ to prove each factor. The input $F_e$ is already
auxiliary independent. Taking only an averaged moment of an unaveraged
$F(R,Y)$ would not suffice: $F=b_e(R)\cos(2\pi Y_1)$ has zero averaged
moment but a nonzero exterior primitive at each generic $Y$.

\section{The source covariance derivative and its signed right inverse}
At each slow point let $w_{\rm tan}=(w_\theta,w_z)$ and define
\[
 \Cov(w)=\av{w_rw_{\rm tan}},\qquad
 \Cross(w,v)=\av{w_rv_{\rm tan}+v_rw_{\rm tan}}.
\]
Direct expansion proves
$\Cov(w+v)=\Cov(w)+\Cross(w,v)+\Cov(v)$ and
$D\Cov(w)[v]=\Cross(w,v)$.

Use the exact source real pulses
$b_\sigma=\chi_g(\xi_g)\psi(v)t_\sigma^h\cos(k\Phi_\sigma)$,
$\sigma\in\{+,-\}$. They exclude the slow cutoff $\eta_\beta$ and
scalar amplitudes; $\chi_g$ is an auxiliary transverse cutoff.
The two signs have disjoint auxiliary supports. Put
$H_{\rm cov}=[\Cov(b_+)\mid\Cov(b_-)]$.
For slow scalar coefficients independent of both averages,
\[
 \Cov(a_+b_++a_-b_-)=H_{\rm cov}(a_+^2,a_-^2)^T.
\]
All cross products vanish by disjointness. The factor $1/2$ from
$\av{\cos^2(k\Phi_\sigma)}_\theta$ is already inside the columns.

On the source interior cone the fixed primary amplitudes satisfy
$a_\sigma=\sqrt{(H_{\rm cov}^{-1}T_{0,*})_\sigma}>0$,
$T_{0,*}=(Q/q)^{A+1/2}T_0$, and
$W_0=\sqrt\varepsilon\sum_\sigma a_\sigma b_\sigma$.
For any signed slow real vector $\Sigma$, define
\begin{equation}
 d_\Sigma=H_{\rm cov}^{-1}\Sigma/\varepsilon,\qquad
 \delta a_\sigma=(d_\Sigma)_\sigma/(2a_\sigma),\qquad
 L_\Sigma=\sqrt\varepsilon\sum_\sigma\delta a_\sigma b_\sigma .
 \label{nsbridge:cb:L}
\end{equation}
It follows exactly that
\begin{gather}
 \Cross(W_0,L_\Sigma)=\varepsilon H_{\rm cov}(2a\odot\delta a)=\Sigma,
 \label{nsbridge:cb:inverse}\\
 \Cov(W_0+L_\Sigma)
   =\varepsilon T_{0,*}+\Sigma+
       \varepsilon H_{\rm cov}(\delta a_+^2,\delta a_-^2)^T.
 \label{nsbridge:cb:quad}
\end{gather}
This proves a linear right inverse of the derivative. It is not an
exact inverse of the quadratic covariance map. There is no sign or
size restriction on $\Sigma$ in these identities; no square root of
$H_{\rm cov}^{-1}(T_{0,*}+\Sigma/\varepsilon)$ is taken.

With Euclidean and induced operator norms and
$a_{\min}=\min_\sigma|a_\sigma|>0$, they give the explicit bound
\begin{equation}
 \|\Cov(L_\Sigma)\|\le
 \frac{\|H_{\rm cov}\|\|H_{\rm cov}^{-1}\|^2}
              {4\varepsilon a_{\min}^2}\,\|\Sigma\|^2. \label{nsbridge:cb:bound}
\end{equation}
Indeed $\|\delta a\|\le\|H_{\rm cov}^{-1}\|\|\Sigma\|/
(2\varepsilon a_{\min})$ and
$\|(\delta a_\sigma^2)_\sigma\|_2\le\|\delta a\|_2^2$.
All higher derivative costs are explicit from Leibniz and
\[
 \partial^I H_{\rm cov}^{-1}
 =-H_{\rm cov}^{-1}
 \sum_{0<J\le I}\binom{I}{J}(\partial^JH_{\rm cov})
                           \partial^{I-J}H_{\rm cov}^{-1},
\]
and, for any nonzero scalar $a$,
$\partial^I(a^{-1})=-a^{-1}\sum_{0<J\le I}\binom{I}{J}
(\partial^Ja)\partial^{I-J}(a^{-1})$.
Together with \eqref{nsbridge:cb:L} these recurrences give finite bounds for
each order without assuming constant amplitudes or constant columns.

\section{Assembly, the actual covariance change, and residual cancellation}
For the physical target $\sigma=(\sigma_2,\sigma_1)$ from
\eqref{nsbridge:cb:primitive}, use the source squared slow partition
$\sum_\beta\eta_\beta^2=1$ on its support. Set
\begin{align}
 W_{0,\rm phys}&=\sum_\beta Q_\beta^{-A}\eta_\beta W_{0,\beta},\nonumber\\
 v_{\rm pr}&=\sum_\beta Q_\beta^{-A}\eta_\beta
                        L_\beta(Q_\beta^{2A}\sigma).
 \label{nsbridge:cb:assembly}
\end{align}
Overlapping slow labels have disjoint auxiliary supports. Taking cross
covariance and using \eqref{nsbridge:cb:inverse} gives
\[
 \Cross(W_{0,\rm phys},v_{\rm pr})
 =\sum_\beta Q_\beta^{-2A}\eta_\beta^2Q_\beta^{2A}\sigma=\sigma.
\]
Every band scale remains distinct. Likewise the leading physical factor is
$Q^{-2A}\varepsilon(Q/q)^{A+1/2}
=Q^{-A+h+1/2}q^{-A-1/2}=q^{-A-1/2}$ because $A=1/2+h$.

Apply the exact curl \eqref{nsbridge:cb:curl} to each coefficient of
$v_{\rm pr}$, including the slow cutoffs before curling. Write the
resulting divergence-free increment as $v_s=v_{\rm pr}+r_s$.
Let $t_{s,m}$ denote the chart transverse harmonic coefficients of
$\eta_\beta L_\beta(Q_\beta^{2A}\sigma)$, before multiplication by
$Q_\beta^{-A}$. For each such coefficient choose
$\pi_{s,m}=i(n^T\mathsf Kt_{s,m}-(n')^Tt_{s,m})/
(k_\beta m|n|^2)$ and sum
$\operatorname{Re}(Q_\beta^{-2A}\pi_{s,m}e^{i\phi_{\beta,m}})$
to define the smooth compact $\pi_s$. Its full gradients remain below.
Let the actual old wave $w$ consist of the wave part of $u_*$
plus the candidate $w_y$, including all existing curl corrections, and write
$w=W_{0,\rm phys}+e_w$. Define tensor covariance
$\mathcal T(w)=\av{w\otimes w}$ and
$\mathcal B_T(u,v)=\av{u\otimes v+v\otimes u}$.
Expansion, with $e_w=w-W_{0,\rm phys}$, proves
\begin{equation}
 \Delta\mathcal T=
 \mathcal B_T(W_{0,\rm phys},v_{\rm pr})
 +\mathcal B_T(e_w,v_{\rm pr})
 +\mathcal B_T(w,r_s)+\mathcal T(v_s).
 \label{nsbridge:cb:fullcov}
\end{equation}
Its first radial--tangential pair is $\sigma$.
All other entries and the remaining three tensors are retained.
For example their radial--tangential pair has bound
\[
 2\|e_w\|_{L^2_{\theta,Y}}\|v_{\rm pr}\|_{L^2_{\theta,Y}}
 +2\|w\|_{L^2_{\theta,Y}}\|r_s\|_{L^2_{\theta,Y}}
 +\|v_s\|_{L^2_{\theta,Y}}^2 ,
\]
by Cauchy--Schwarz. The compact-set estimates above and the explicit
curl differentiate to bound all these terms at each finite order.

For clarity the complete residual after adding $v_s$ and its selected
pressure $\pi_s$ is, without any averaging,
\begin{align}
 R_{\nu_{\rm NS}}(u+v_s,p+\pi_s)
 =R_{\nu_{\rm NS}}(u,p)
 &+\partial_tv_s+(u\cdot\nabla)v_s+(v_s\cdot\nabla)u \nonumber\\
 &-\nu_{\rm NS}\Delta v_s+\nabla\pi_s
 +(v_s\cdot\nabla)v_s . \label{nsbridge:cb:residual}
\end{align}
This follows by expanding the product and retains viscosity and every
pressure term. For a symmetric tensor $S$ the averaged tangential
divergence in physical cylindrical coordinates is
\[
 (\av{\nabla\cdot S})_\theta
   =(\partial_r+2/r)\av{S_{r\theta}}+\partial_z\av{S_{z\theta}},
 \qquad
 (\av{\nabla\cdot S})_z
   =(\partial_r+1/r)\av{S_{rz}}+\partial_z\av{S_{zz}} .
\]
To verify, write the Cartesian divergence in the orthonormal cylindrical
basis, using $\partial_\theta e_r=e_\theta$,
$\partial_\theta e_\theta=-e_r$; periodic angular derivatives average
to zero and the two basis terms give $2/r$ and $1/r$ respectively.
The auxiliary chain terms average to zero by periodicity, but the
physical slow derivatives remain.

Let $U=u_*+w_y$, $P_y=p_*+\pi_y$, and set
\[
 \mathcal L=
 \partial_tv_s+(U\cdot\nabla)v_s+(v_s\cdot\nabla)U
 -\nu_{\rm NS}\Delta v_s+\nabla\pi_s .
\]
For a fully explicit averaged remainder, write $U=u_{\rm mean}+w$,
where $u_{\rm mean}$ is its angular mean, and put
$\mathcal V=\mathcal L-(w\cdot\nabla)v_s-(v_s\cdot\nabla)w$.
Then the average of \eqref{nsbridge:cb:residual} is
$\av{R_y}+\av{\mathcal V}+\nabla\cdot\Delta\mathcal T$.
Let $S_0=\mathcal B_T(W_{0,\rm phys},v_{\rm pr})$ and
$S_{\rm rem}=\mathcal B_T(e_w,v_{\rm pr})
 +\mathcal B_T(w,r_s)+\mathcal T(v_s)$.
Consequently the radial divergence of the first pair in
\eqref{nsbridge:cb:fullcov} changes the prescribed averaged residual $F_e$ into
$b_eM_e$ by \eqref{nsbridge:cb:radial}. The axial tensor divergences, the other
three tensors, the exact linear terms and the pressure terms of
\eqref{nsbridge:cb:residual} remain in the new residual. Equivalently the new
averaged tangential residual is exactly
\begin{align*}
 \av{R_{\rm new}}_{\rm tan}
 ={}&\av{R_{\nu_{\rm NS}}(u_*,p_*)}_{\rm tan}
 +(b_2M_2,b_1M_1)+\av{\mathcal V}_{\rm tan}\\
 &+\partial_z((S_0)_{z\theta},(S_0)_{zz})
   +(\nabla\cdot S_{\rm rem})_{\rm tan}.
\end{align*}
Every term here is defined by the displayed differentiated fields and
tensors. None is assumed to vanish.
The map from $y$ through \eqref{nsbridge:cb:F}, \eqref{nsbridge:cb:primitive},
\eqref{nsbridge:cb:L}, and \eqref{nsbridge:cb:assembly} is now an actual composite
construction. It retains the signed residual produced by the coupled
candidate and the source's original physical and chart units.

\section{Signs and the exact full-flow reflection}
The maps $F_e\mapsto M_e\mapsto\sigma_e\mapsto L_\Sigma$ are linear
and reverse sign with their inputs; $\Cov(L_\Sigma)$ is even.
The two primary amplitudes may also be replaced independently by
$s_\sigma a_\sigma$, $s_\sigma\in\{-1,1\}$, provided the same fixed
choice is used in each denominator: the cross identity is unchanged.
This does not negate the nonlinear Navier--Stokes equation.

Here is the exact full-flow map. In Cartesian coordinates let
$O=\operatorname{diag}(1,-1,1)$ and define
\[
 u^\sharp(x,t)=O u(Ox,t),\quad
 p^\sharp(x,t)=p(Ox,t),\quad f^\sharp(x,t)=O f(Ox,t).
\]
The chain rule gives
$\nabla u^\sharp=O(\nabla u)(Ox,t)O$,
$\Delta u^\sharp=O(\Delta u)(Ox,t)$, and
$(u^\sharp\cdot\nabla)u^\sharp
 =O((u\cdot\nabla)u)(Ox,t)$.
Thus $\nabla\cdot u^\sharp=0$ and
$R_{\nu_{\rm NS}}(u^\sharp,p^\sharp)=O R_{\nu_{\rm NS}}(u,p)(Ox,t)$
with the same positive viscosity.
In cylindrical components this is
\[
 (u_r^\sharp,u_\theta^\sharp,u_z^\sharp)(r,\theta,z,t)
 =(u_r,-u_\theta,u_z)(r,-\theta,z,t),
\]
and the force transforms identically. Pressure reflects in its argument.
For an axisymmetric field this reverses exactly its swirl; meridional
components remain fixed. The entire momentum tensor is mapped to
$O\mathcal T O^T$, so its radial--tangential pair transforms by
$(T_{r\theta},T_{rz})\mapsto(-T_{r\theta},T_{rz})$ after angular
averaging. Vorticity transforms as
$\omega^\sharp(x,t)=\det(O)\,O\omega(Ox,t)$, obtained by applying
the same derivative identities to the Levi-Civita formula for curl.
The map is an involution, preserves the kinetic-energy integral
(by $|\det O|=1$), compact support, smoothness, and blowup magnitude
along the reflected path. It preserves zero initial data when present.
It proves transport of any established axisymmetric signed swirl
divergence to the opposite sign; it does not itself prove an endpoint.

The source five-moment correction, auxiliary mean inverse, and infinite
summation are subsequent source operations. Their full estimates are
not re-proved here. The earlier draft formulas for them are withdrawn:
in particular the fifth moment row is
$\int(2RG\gamma_d-RV\Delta v)\,\dd R=-J_z$, without an extra outer
$R$, and its reserved profile in a $Q$ chart includes $(Q/q)^A$.
The finite map proved above retains the two moments and all other
residual terms, so it establishes neither a third shooting return nor
an infinite modified viscous sequence.


\section{The released result and its actual signed path}
\label{nsprop:source}
The source in this section is the supplied 166-page edition of
\emph{Finite Time Blowup for Navier--Stokes}, whose SHA-256 is
\begin{center}\footnotesize
\texttt{0e779481c4da40bd28d1e642e1d8ca57447d129610df28dfa5a11e9af8ae228f}.
\end{center}
Its Theorem 1.1, on page 1, states the following result for the standard
three-dimensional incompressible Navier--Stokes equation. For every
fixed physical viscosity $\nu_{\rm NS}>0$, the source constructs
\[
 f\in C^\infty_c(\mathbb R^3\times(0,\infty);\mathbb R^3),
 \qquad u,p\in C^\infty(\mathbb R^3\times[0,1)),
\]
with a fixed compact spatial support $K$ for $u$ and $p$, such that
\begin{equation}
 \begin{gathered}
 \partial_tu+(u\cdot\nabla)u-\nu_{\rm NS}\Delta u+\nabla p=f,
 \quad \nabla\cdot u=0,\quad u(\cdot,0)=0,\\
 \quad \sup_{t<1}\|u(t)\|_2<\infty,\quad
 \limsup_{t\uparrow1}\|u(t)\|_\infty=\infty .
 \end{gathered}
 \label{nsprop:source_theorem}
\end{equation}
It also excludes a global smooth solution with the same force and datum
and uniformly bounded kinetic energy. Its Corollary 10.6, pages 125--126,
gives a periodic construction on $\mathbb R^3/\mathbb Z^3$ with smooth
time-compact force, zero initial datum and no global smooth continuation.
The source identifies these with alternatives (C) and (D) in the
referenced Millennium problem statement. The force is nonzero, as the
source explicitly notes on page 124. These are statements of the cited
source theorem; the finite algebra checks in this lane do not constitute
an independent verification of the complete 166-page proof.

The source proves more than an unsigned norm divergence. In its
viscosity-one construction, equations (10.20)--(10.21), page 124, give
\begin{equation}
 \begin{split}
 A&=\tfrac12+h,\qquad 0<h<1/100,\qquad X_{\rm in}>0,\quad e_0>0,\\
 x_\tau&=(\sqrt{2X_{\rm in}\tau},0,0),\qquad t_\tau=1-\tau,\\
 u_\theta(x_\tau,t_\tau)
   &=\tau^{-A}(e_0+\epsilon(\tau)),\qquad
      |\epsilon(\tau)|\le C\tau^{2h}\quad(0<\tau<\tau_0).
 \end{split}\label{nsprop:positive_path}
\end{equation}
Here $C,\tau_0$ are finite constants supplied by the source asymptotic;
$e_0$ and $X_{\rm in}$ retain their source values. If $C>0$, choose
$0<\tau_1\le\tau_0$ so that $C\tau_1^{2h}\le e_0/2$; if $C=0$,
take $\tau_1=\tau_0$. It follows directly that
\[
 u_\theta(x_\tau,t_\tau)\ge \frac{e_0}{2}\tau^{-A}
       \longrightarrow+\infty\qquad(0<\tau<\tau_1).
\]
The observation point moves to the origin. This assertion does not
replace that path by the fixed point $x=0$. On the path, $q=\tau>0$,
and at its cylindrical angle $\theta=0$ one has $e_\theta=e_2$.
The geometric domain $q>0$, the positive exponent $h$, and the signed
velocity value are three retained parts of the same explicit formula.

\section{Exact propagation at retained viscosities and orientations}
\label{nsprop:map}
For any smooth fields $u,p$ define the physical momentum residual
\[
 R_{\nu_b}(u,p)=\partial_tu+(u\cdot\nabla)u-\nu_b\Delta u+\nabla p.
\]
Let $\nu_b>0$ and $\nu_t>0$ be the base and target physical viscosities,
and retain
\[
 \rho=\sqrt{\nu_t/\nu_b},\qquad O^TO=OO^T=I,\qquad a\in\mathbb R^3.
\]
Define the affine spatial coordinate $y=O^T(x-a)/\rho$ and the fields
\begin{equation}
 u^\sharp(x,t)=\rho O u(y,t),\quad
 p^\sharp(x,t)=\rho^2p(y,t),\quad
 f^\sharp(x,t)=\rho O f(y,t).
 \label{nsprop:affine_map}
\end{equation}
No time coordinate is changed. These formulas give a bijection on smooth
field triples, with inverse
\[
 u(y,t)=\rho^{-1}O^Tu^\sharp(a+\rho Oy,t),\quad
 p(y,t)=\rho^{-2}p^\sharp(a+\rho Oy,t),\quad
 f(y,t)=\rho^{-1}O^Tf^\sharp(a+\rho Oy,t).
\]
For the Jacobian convention $(Du)_{ij}=\partial_j u_i$, differentiation
gives, without dropping an orientation or viscosity,
\begin{align*}
 Du^\sharp&=O(Du)(y,t)O^T,&
 \nabla\cdot u^\sharp&=(\nabla\cdot u)(y,t),\\
 \partial_tu^\sharp&=\rho O\partial_tu(y,t),&
 (u^\sharp\cdot\nabla_x)u^\sharp&=\rho O((u\cdot\nabla_y)u)(y,t),\\
 \Delta_xu^\sharp&=\rho^{-1}O\Delta_yu(y,t),&
 \nabla_xp^\sharp&=\rho O\nabla_yp(y,t).
\end{align*}
The nonlinear identity follows by multiplying $Du^\sharp$ by
$u^\sharp$ and using $O^TO=I$. The Laplacian identity follows by
contracting the two spatial derivative indices and using orthogonality.
Since $\nu_t=\rho^2\nu_b$, these identities prove
\begin{equation}
 R_{\nu_t}(u^\sharp,p^\sharp)-f^\sharp
       =\rho O\bigl(R_{\nu_b}(u,p)-f\bigr)(y,t).
 \label{nsprop:residual_map}
\end{equation}
In particular the map preserves the exact equation as well as each
nonzero residual with its exact factor.

For vorticity, the cross-product identity
$(Ov)\times(Ow)=\det(O)O(v\times w)$ gives
\begin{equation}
 \omega^\sharp(x,t)=\det(O)O\omega(y,t).
 \label{nsprop:vorticity}
\end{equation}
One proof of this identity takes its scalar product with $Oz$:
both sides give $\det(O)\det[v,w,z]$ for every $z$. Applying it to
the contracted first derivatives proves the vorticity formula; the
factor $\rho$ in velocity cancels the inverse factor in a derivative.

Spatial support becomes $a+\rho OK$. The Jacobian has absolute
determinant $\rho^3$, for both signs of $\det O$. Hence, for
$1\le p<\infty$,
\begin{equation}
 \|u^\sharp(t)\|_p=\rho^{1+3/p}\|u(t)\|_p,\qquad
 \|u^\sharp(t)\|_\infty=\rho\|u(t)\|_\infty,\qquad
 \|u^\sharp(t)\|_2^2=\rho^5\|u(t)\|_2^2 .
 \label{nsprop:norms}
\end{equation}
The Frobenius norm of $O(Du)O^T$ equals that of $Du$. Consequently,
with integrals over the same time interval,
\[
 \nu_t\int\|\nabla u^\sharp\|_2^2\,dt
     =\rho^5\nu_b\int\|\nabla u\|_2^2\,dt .
\]
For arbitrary directions $v_1,\ldots,v_j$, the exact force derivative is
\[
 D_x^j\partial_t^m f^\sharp[v_1,\ldots,v_j](x,t)
 =\rho^{1-j}O
   (D_y^j\partial_t^m f)[O^Tv_1,\ldots,O^Tv_j](y,t).
\]
Thus smoothness, compact space-time support and vanishing of every
derivative at a terminal point are preserved. The initial datum remains
zero. A global smooth bounded-energy competitor for the transformed
datum and force would, under the displayed inverse, be a competitor
for the base datum and force, because its energy is multiplied by the
fixed factor $\rho^{-5}$. This proves propagation of the source
noncontinuation conclusion through the map.

\section{The source positive branch and the exact negative branch}
\label{nsprop:branches}
Apply the map to the source field in \eqref{nsprop:positive_path}.
Here $\nu_b=1$ specifies the source object being used; physical target
viscosity $\nu_t$ remains arbitrary and $\rho=\sqrt{\nu_t}$ is retained.
Take $a=0$ and the two particular orthogonal matrices
\[
 O_+=I,\qquad O_-=\operatorname{diag}(1,-1,1).
\]
Both fix $e_1$, so they give the same spatial observation path
$x^\sharp_\tau=\rho x_\tau$ and the same time $1-\tau$.
Since $O_+e_2=e_2$ and $O_-e_2=-e_2$, their exact signed values are
\begin{equation}
 u^\sharp_{\pm,\theta}(\rho x_\tau,1-\tau)
     =\pm\rho\,\tau^{-A}(e_0+\epsilon(\tau))
        \longrightarrow\begin{cases}+\infty,&+,\\-\infty,&-.\end{cases}
 \label{nsprop:two_paths}
\end{equation}
All support, energy, force and viscosity statements in the preceding
section apply to both triples. For the negative branch the force is
$\rho O_-f(O_-x/\rho,t)$. Equality of this force with the positive
branch force is not needed and is not asserted.

For a general, possibly nonaxisymmetric field the reflection alone has
the exact cylindrical component formula
\[
 (u^r,u^\theta,u^z)^\sharp(r,\theta,z,t)
    =(u^r,-u^\theta,u^z)(r,-\theta,z,t).
\]
Indeed $O_-e_r(-\theta)=e_r(\theta)$,
$O_-e_\theta(-\theta)=-e_\theta(\theta)$ and $O_-e_z=e_z$.
This proves the formula even in the presence of all oscillatory
corrections. For the axisymmetric background the argument $-\theta$
can be dropped because its components are independent of $\theta$.
For a velocity covariance tensor the exact Cartesian transformation is
$T^\sharp(x,t)=\rho^2 O T(y,t)O^T$: expand the outer product of
$\rho Ow$ with itself and use linearity of the retained average.
For its angularly averaged symmetric tensor, the reflection basis
calculation at $\rho=1$ gives
\[
 (T_{r\theta},T_{rz})^\sharp=(-T_{r\theta},T_{rz}).
\]
Under the additional viscosity dilation its components have factor
$\rho^2$ at the corresponding points. Thus the two covariance columns
and their target transform together. Positive coefficients in a source
representation $T=c_1v_1+c_2v_2$, $c_i>0$, remain the same coefficients
of the transformed columns; no change of coefficient sign is required
to realize the negative-swirl branch.

The source's periodic construction also has this reflection: $O_-$
preserves $\mathbb Z^3$, so $[x]\mapsto[O_-x]$ is well defined on
$\mathbb R^3/\mathbb Z^3$. The derivative identities apply to periodic
lifts. This gives both signed orientations of the source Corollary 10.6.
Arbitrary $O\in O(3)$ need not preserve this particular lattice; only
the lattice-preserving reflection just specified is used for this
periodic assertion.

\section{Propagation into the coupled calculation}
\label{nsprop:programme}
The preceding sections use the released whole-space and periodic
constructions at the precise places where those source results apply.
The original coupled Boussinesq work is due to Levent Alp\"oge and
Tristan Buckmaster. This lane's retained coupled system, its finite
controlled stages and its residual maps remain the separate objects
defined in the preceding parts of this reader.

For every compact candidate $(u_*+w_y,p_*+\pi_y)$ in the covariance
part, \eqref{nsprop:residual_map} propagates its full computed residual.
Linearity of the transformation and the exact nonlinear identity give
in particular
\[
 \Delta R_y^\sharp(x,t)=\rho O\Delta R_y(y,t).
\]
Thus the finite candidate's cutoff forces, curl errors and covariance
remainders all have an exact transformed counterpart. A nonzero
remainder is carried to a nonzero remainder, because the map is
invertible. The source's residual flatness at the singular point belongs
to its completed local source construction; its extended smooth force
need not vanish on the entire time-one slice. The next section calculates the
surviving moments of the actual finite candidate, so that propagation
continues through its computed terms rather than assuming that its
different coupled coefficients inherit the source's infinite estimates.

The earlier suggestion of a negative blowout imposes no condition on
the programme. Formula \eqref{nsprop:positive_path} is the source-selected
sign; formula \eqref{nsprop:two_paths} proves the exact relation of its
two orientations. Both formulas retain a positive physical viscosity
and the source coordinate domain $q>0$.


\section{The two moments of the actual compact curl increment}
\label{nsmom:setup}
The objects in this calculation are the finite smooth source state
$(u_*,p_*)$ and the actual added curl field $v=w_y$ with its selected
pressure $\pi=\pi_y$ from the signed covariance bridge. Their physical
residual increment is
\begin{equation}
 \Delta R=\partial_t v+(u_*\cdot\nabla)v+(v\cdot\nabla)u_*
 +(v\cdot\nabla)v-\nu_{\rm NS}\Delta v+\nabla\pi .
 \label{nsmom:increment}
\end{equation}
The physical viscosity $\nu_{\rm NS}$ is unchanged. Both velocities are
divergence free: the source state uses its actual divergence-free
background and curl increments, and $v$ is the curl of its specified
compact potential. No residual equation for $u_*$ is required in
\eqref{nsmom:increment}; its residual has been subtracted exactly.
All radial supports involving $v$ or $\pi$ lie in a fixed compact
subinterval of $(0,\infty)$ on the compact space--time patch at issue.

Here the source extended field is indispensable. Its absolute auxiliary
coordinate is $Y\in\mathbb T^2=\mathbb R^2/\mathbb Z^2$, and physical
evaluation uses precisely source (6.2)--(6.4):
\begin{gather}
 J_g=\begin{pmatrix}3&1\\1&5\end{pmatrix},\quad
 \Lambda_g=4-\sqrt2,\quad T_g=4+\sqrt2,\quad b_g=\sqrt2-1,\nonumber\\
 \mathbf v_r=(1,-b_g),\quad\mathbf v_t=(b_g,1),\quad
 \rho_g=\frac{\log\Lambda_g}{\log T_g},\quad
 \kappa_s=10^{-5},\quad d_r=2((1+h)\rho_g-h\kappa_s),\nonumber\\
 Y(r,t)=\mathbf v_r r^{d_r}+\mathbf v_t t\pmod{\mathbb Z^2},\qquad
 \mathscr D_r=\partial_r+d_r r^{d_r-1}\mathbf v_r\cdot\partial_Y,
 \quad\mathscr D_t=\partial_t+\mathbf v_t\cdot\partial_Y.
 \label{nsmom:phase}
\end{gather}
These are the source's original quantities, not the coupled physical
frequency or diffusivity. Angular and axial physical derivatives are
$\partial_\theta$ and $\partial_z$, because $Y(r,t)$ depends on neither
variable. Every field below is first differentiated using these lifted
physical derivatives. A tilde is suppressed for this extended
representation. Evaluation at \eqref{nsmom:phase} recovers the physical
field and its physical derivatives by the chain rule.

The incompressibility needed below holds at every independent $Y$.
Indeed the actual finite source construction retains the form
$u_* =\operatorname{curl}_{\mathscr D}\mathcal A_*+
B_*e_\theta$ with $\partial_\theta B_*=0$, and the added field is
$v=\operatorname{curl}_{\mathscr D}\mathcal A_y$. The retained source
potentials here include its finite mean and wave curl increments.
The lifted derivatives in \eqref{nsmom:phase}, together with
$\partial_\theta$ and $\partial_z$, commute, and
$\mathscr D_r r=1$. Substitution in the cylindrical curl and
divergence formulas therefore gives
$\operatorname{div}_{\mathscr D}\operatorname{curl}_{\mathscr D}
\mathcal A=0$ identically in $(r,\theta,z,t,Y)$: each mixed derivative
cancels its opposite, including the derivative of $r^{-1}$.
Also $\operatorname{div}_{\mathscr D}(B_*e_\theta)
=r^{-1}\partial_\theta B_*=0$.
Thus extended incompressibility is proved from the actual retained
potential representation, rather than inferred from its restriction to
the physical phase map.

Use the normalized averages
\[
 \langle a\rangle_\theta=\frac1{2\pi}\int_0^{2\pi}a\,d\theta,
 \qquad \langle a\rangle_Y=\int_{\mathbb T^2}a\,dY,
 \qquad\langle a\rangle=\langle\langle a\rangle_\theta\rangle_Y,
 \qquad\int_{\mathbb T^2}1\,dY=1.
\]
They act on cylindrical components before physical evaluation, exactly
as source (3.7)--(3.9). All claimed integral identities below concern
these specified averages; the exact relation to evaluated moments is
also proved below.

\subsection{The actual harmonic curl has zero angular mean}
\label{nsmom:zero_mean}
For each retained source harmonic, its extended potential and pressure
have cylindrical components of the form
\[
 C_j(r,z,t,Y)e^{i(n_j\theta+\psi_j(r,z,t,Y))},\qquad
 P_j(r,z,t,Y)e^{i(n_j\theta+\psi_j(r,z,t,Y))},\qquad
 n_j=k_{\beta_j}m_jp_j\in\mathbb Z\setminus\{0\}.
\]
The coefficients include the actual cutoff $\chi$, frame $B$, map $L$,
coupled solution $y=(\Theta,\Omega)^T$, band powers and pressure formula.
Their cylindrical components are independent of $\theta$. In particular,
no sign of $\Theta$, $\Omega$ or the real part of a harmonic is prescribed.
For a vector potential $\mathcal A$ the exact cylindrical curl is
\[
 (\operatorname{curl}\mathcal A)_r=r^{-1}\partial_\theta\mathcal A_z
                         -\partial_z\mathcal A_\theta,
 \quad
 (\operatorname{curl}\mathcal A)_\theta=\partial_z\mathcal A_r
                         -\mathscr D_r\mathcal A_z,
 \quad
 (\operatorname{curl}\mathcal A)_z=(\mathscr D_r+r^{-1})\mathcal A_\theta
                         -r^{-1}\partial_\theta\mathcal A_r.
\]
These formulas include the derivatives of the cylindrical basis.
Each displayed operation retains the same factor $e^{in_j\theta}$:
the coefficients of $\mathscr D_r$ are independent of $\theta$, and
$\partial_\theta$ multiplies that factor by $in_j$. Integration of the
factor over the angular circle gives zero. Taking real parts and the
actual finite harmonic sum proves
\begin{equation}
 \langle v_r\rangle_\theta=\langle v_\theta\rangle_\theta
 =\langle v_z\rangle_\theta=\langle\pi\rangle_\theta=0
 \quad\hbox{at every }(r,z,t,Y).
 \label{nsmom:mean_zero}
\end{equation}
This includes every curl and cutoff correction, not merely the
principal amplitude. The same identities hold after physical
evaluation because $Y(r,t)$ is independent of $\theta$.

For every smooth periodic coefficient $a$, integration of its auxiliary
derivatives gives
\begin{equation}
 \langle\mathscr D_r a\rangle=\partial_r\langle a\rangle,
 \quad\langle\mathscr D_t a\rangle=\partial_t\langle a\rangle,
 \quad\langle\partial_z a\rangle=\partial_z\langle a\rangle,
 \quad\langle\partial_\theta a\rangle=0.
 \label{nsmom:commute}
\end{equation}
There is no omitted variable coefficient in this assertion:
$d_r r^{d_r-1}\mathbf v_r$ is independent of $Y$. Applying the first
identity again proves
$\langle\mathscr D_r^2a\rangle=\partial_r^2\langle a\rangle$.
Explicitly, if $c(r)=d_r r^{d_r-1}\mathbf v_r$, then
\[
 \mathscr D_r^2 a=\partial_r^2a+2c\cdot\partial_Y\partial_r a
 +c'\cdot\partial_Ya+(c\cdot\partial_Y)^2a;
\]
each of the three additional terms has zero auxiliary average.

\subsection{Exact removal of the linear and mean-flow terms}
\label{nsmom:linear}
Put
\[
 \bar u_*:=\langle u_*\rangle_\theta,
 \qquad w_*:=u_*-\bar u_*.
\]
The field $\bar u_*$ may retain auxiliary dependence; none is discarded.
Since it is independent of $\theta$, \eqref{nsmom:mean_zero} gives
\[
 \langle\bar u_*\otimes v+v\otimes\bar u_*\rangle=0.
\]
Define the actual symmetric cross-and-quadratic tensor
\begin{equation}
 S:=\langle w_*\otimes v+v\otimes w_*+v\otimes v\rangle.
 \label{nsmom:stress}
\end{equation}
Products of all original source waves with the added curl field are
retained in this formula, including old and new curl remainders.

For two divergence-free fields $a,b$, Cartesian component expansion
gives
\[
 \nabla\cdot(a\otimes b+b\otimes a)
 =(b\cdot\nabla)a+(a\cdot\nabla)b,
 \qquad\nabla\cdot(v\otimes v)=(v\cdot\nabla)v.
\]
For example
$\sum_j\partial_j(a_i b_j)=\sum_jb_j\partial_ja_i+
a_i\sum_j\partial_jb_j$; the last term is zero. The lifted derivatives
obey the same product rule, so the identities hold on the extended
domain before evaluation. The three advection terms in
\eqref{nsmom:increment} are therefore exactly the divergence of
$u_*\otimes v+v\otimes u_*+v\otimes v$.

For completeness, set
$\mathscr L_0=\mathscr D_r^2+r^{-1}\mathscr D_r+
r^{-2}\partial_\theta^2+\partial_z^2$.
The vector Laplacian and pressure gradient in the two components are
\[
 (\Delta v)_\theta=(\mathscr L_0-r^{-2})v_\theta
                       +2r^{-2}\partial_\theta v_r,
 \qquad(\Delta v)_z=\mathscr L_0v_z,
 \qquad(\nabla\pi)_\theta=r^{-1}\partial_\theta\pi,
 \qquad(\nabla\pi)_z=\partial_z\pi.
\]
Equations \eqref{nsmom:mean_zero}--\eqref{nsmom:commute} prove separately
\begin{equation}
 \langle\mathscr D_t v\rangle_{\rm tan}=0,
 \qquad -\nu_{\rm NS}\langle\Delta v\rangle_{\rm tan}=0,
 \qquad\langle\nabla\pi\rangle_{\rm tan}=0.
 \label{nsmom:linear_zero}
\end{equation}
Thus viscosity and pressure were included before averaging; their
vanishing here is a proved consequence of the actual carriers.

If $T$ is symmetric, differentiation of
$e_r=(\cos\theta,\sin\theta,0)$,
$e_\theta=(-\sin\theta,\cos\theta,0)$ gives
\begin{align*}
 (\nabla\cdot T)_\theta
   &=(\mathscr D_r+2/r)T_{r\theta}
       +r^{-1}\partial_\theta T_{\theta\theta}+\partial_zT_{z\theta},\\
 (\nabla\cdot T)_z
   &=(\mathscr D_r+1/r)T_{rz}
       +r^{-1}\partial_\theta T_{\theta z}+\partial_zT_{zz}.
\end{align*}
Indeed $\partial_\theta e_r=e_\theta$ and
$\partial_\theta e_\theta=-e_r$ supply the two radial connection
terms in the first line and the single radial connection term in the
second. Averaging these formulas and using
\eqref{nsmom:commute}, \eqref{nsmom:stress}, and \eqref{nsmom:linear_zero}
proves the exact residual identities
\begin{equation}
 F_2:=\langle\Delta R_\theta\rangle
    =(\partial_r+2/r)S_{r\theta}+\partial_zS_{z\theta},
 \qquad
 F_1:=\langle\Delta R_z\rangle
    =(\partial_r+1/r)S_{rz}+\partial_zS_{zz}.
 \label{nsmom:F}
\end{equation}
In particular every mean-flow cross term vanishes as the divergence
of its zero averaged tensor, even when the mean depends on $Y$.

\subsection{The surviving axial fluxes and their signs}
\label{nsmom:moments}
Define the two moments with the original physical radial measure:
\begin{equation}
 M_2(z,t)=\int_0^\infty r^2F_2(r,z,t)\,dr,
 \qquad M_1(z,t)=\int_0^\infty rF_1(r,z,t)\,dr.
 \label{nsmom:Mdef}
\end{equation}
Their exact values are
\begin{align}
 M_2&=\partial_z J_2, &
 J_2&=\int_0^\infty r^2
    \langle w_{*,z}v_\theta+v_z w_{*,\theta}+v_zv_\theta\rangle\,dr,
 \label{nsmom:M2}\\
 M_1&=\partial_z J_1, &
 J_1&=\int_0^\infty r
    \langle2w_{*,z}v_z+v_z^2\rangle\,dr.
 \label{nsmom:M1}
\end{align}
To prove them, multiply \eqref{nsmom:F} by the respective weight:
$r^2(\partial_r+2/r)S_{r\theta}=\partial_r(r^2S_{r\theta})$
and $r(\partial_r+1/r)S_{rz}=\partial_r(rS_{rz})$.
The weighted boundary values at both ends of $(0,\infty)$ are zero,
because each entry of $S$ contains at least one compactly supported
factor $v$. Differentiation under the radial integral is valid by the
same compact support and smoothness. This proves
\eqref{nsmom:M2}--\eqref{nsmom:M1} including their signs and factors.

There is no pointwise cancellation of the axial derivatives in these
formulas. Since the added field has compact axial support, $J_1$ and
$J_2$ have compact axial support and
\begin{equation}
 \int_{\mathbb R}M_e(z,t)\,dz=0\quad(e=1,2).
 \label{nsmom:global_zero}
\end{equation}
This follows by evaluating the boundary values of $J_e$ at the two
axial ends. It does not assert $M_e(z,t)=0$ at an individual axial
point. The compact radial primitive in the covariance bridge must
therefore retain its $b_eM_e$ terms with these actual $M_e$.

At fixed source state, replacing the added field and pressure by their
negatives sends the old-wave cross terms in \eqref{nsmom:M2}--\eqref{nsmom:M1}
to their negatives and leaves the quadratic terms unchanged. Thus
negating an amplitude is not a sign reversal of either full moment.
Under the full physical reflection
$O=\operatorname{diag}(1,-1,1)$, applied to the source state and added
field together as $a^\sharp(x,t)=Oa(Ox,t)$, the cylindrical radial and
axial components retain their signs and the angular components reverse
sign at $(r,-\theta,z,t)$. Changing the angular integration variable
therefore proves
\[
 J_2^\sharp=-J_2,\quad M_2^\sharp=-M_2,
 \qquad J_1^\sharp=J_1,\quad M_1^\sharp=M_1.
\]

An explicit transverse-curl example shows why compactness and nonzero
angular frequency alone do not force either moment to vanish. This
example tests precisely those structural properties, and is not a
replacement for the retained source state. Let $H$ be any nonzero real
$C_c^\infty((r_-,r_+))$ function with $0<r_-<r_+$, let
$g\in C_c^\infty(\mathbb R)$ be nonconstant, let $n$ be a nonzero
integer, and put $f=rH'$. Use zero old state and the real potential
\[
 \mathcal A_r=fg\sin(n\theta),\qquad
 \mathcal A_\theta=0,\qquad
 \mathcal A_z=Hg\cos(n\theta).
\]
It has the required transverse representation:
$\phi=n\theta$, $\xi=(n/r)e_\theta$,
$C=(-ifg,0,Hg)$ in cylindrical components,
$a=i\xi\times C$, and
$C=i\xi\times a/|\xi|^2$ because $\xi\cdot C=0$.
The curl is exactly
\[
 v_r=-\frac{nHg}{r}\sin(n\theta),\qquad
 v_\theta=fg'\sin(n\theta)-H'g\cos(n\theta),\qquad
 v_z=-\frac{nfg}{r}\cos(n\theta).
\]
Consequently
\[
 M_2=ngg'\int_0^\infty r^2(H')^2\,dr,
 \qquad
 M_1=n^2gg'\int_0^\infty r(H')^2\,dr.
\]
For example $g(z)=\exp(-1/(1-z^2))$ on $|z|<1$, extended by zero,
has $g(z)g'(z)\ne0$ at every $0<|z|<1$. Both integrals are strictly
positive: a compactly supported smooth function with $H'=0$ everywhere
would be constant and hence zero. Thus both moments can be nonzero,
with their displayed signed values, within the exact transverse-curl
architecture.

\subsection{Bounds in the actual fields and original chart units}
\label{nsmom:bounds}
At fixed $(z,t)$ set
\[
 \|a\|_e^2=\int_0^\infty r^e\langle|a|^2\rangle\,dr,
 \qquad e=1,2.
\]
When an old-field norm is used below it can be restricted to the
enclosing radial support of $v$ and $\partial_zv$. This avoids imposing
any decay condition on an unused exterior part of the old field.
The following component bounds contain only actual fields:
\begin{align}
 |J_2|\le{}&\|w_{*,z}\|_2\|v_\theta\|_2
 +\|w_{*,\theta}\|_2\|v_z\|_2+\|v_z\|_2\|v_\theta\|_2,
 \label{nsmom:J2bound}\\
 |J_1|\le{}&2\|w_{*,z}\|_1\|v_z\|_1+\|v_z\|_1^2,
 \label{nsmom:J1bound}\\
 |M_2|\le{}&\|\partial_z w_{*,z}\|_2\|v_\theta\|_2
 +\|w_{*,z}\|_2\|\partial_zv_\theta\|_2
 +\|\partial_zv_z\|_2\|w_{*,\theta}\|_2\nonumber\\
 &+\|v_z\|_2\|\partial_z w_{*,\theta}\|_2
 +\|\partial_zv_z\|_2\|v_\theta\|_2
 +\|v_z\|_2\|\partial_zv_\theta\|_2,
 \label{nsmom:M2bound}\\
 |M_1|\le{}&2\bigl(\|\partial_z w_{*,z}\|_1\|v_z\|_1
 +\|w_{*,z}\|_1\|\partial_zv_z\|_1
 +\|v_z\|_1\|\partial_zv_z\|_1\bigr).
 \label{nsmom:M1bound}
\end{align}
Each is the Cauchy--Schwarz inequality on the product measure
$r^e\,dr\,d\theta/(2\pi)\,dY$, after applying the product rule in
\eqref{nsmom:M2}--\eqref{nsmom:M1}. For example, the derivative of
$2w_{*,z}v_z+v_z^2$ is
$2(\partial_z w_{*,z})v_z+2w_{*,z}\partial_zv_z+
2v_z\partial_zv_z$, proving every coefficient in
\eqref{nsmom:M1bound}.
For a $(z,t)$ multi-index $I$, every further derivative is exactly
obtained by replacing each product $ab$ in $J_e$ by
\[
 \partial^{I+(1,0)}(ab)
 =\sum_{K\le I+(1,0)}\binom{I+(1,0)}{K}
       (\partial^Ka)(\partial^{I+(1,0)-K}b)
\]
inside its original radial integral. This identity follows inductively
from the one-derivative product rule, and gives the corresponding
finite-order moment bound by applying the same Cauchy--Schwarz
inequality separately to every displayed summand.

The field $v$ in these bounds is not an abstract replacement amplitude.
For each actual band and harmonic write
\begin{gather}
 t_j=\chi_jB_jL_jy_j,\qquad
 c_j=\frac{i\,n_{\Phi,j}\times t_j}
                  {k_{\beta_j}m_j|n_{\Phi,j}|^2},\qquad
 b_j=t_j+\operatorname{curl}_{*,\mathrm{coeff}}c_j,\nonumber\\
 v=\sum_j\operatorname{Re}
       \bigl(Q_{\beta_j}^{-A}b_j e^{ik_{\beta_j}m_j\Phi_j}\bigr),
 \qquad A=\tfrac12+h,\quad D=\tfrac12-h.
 \label{nsmom:actual_amplitude}
\end{gather}
Here $\operatorname{curl}_{*,\mathrm{coeff}}$ differentiates the chart
coefficient, including the cutoff, auxiliary chain rule and cylindrical
basis, but not the displayed exponential. This is exactly the
remainder of the physical potential
$Q_{\beta_j}^{1/2-A}c_j e^{ik_{\beta_j}m_j\Phi_j}$: physical curl
introduces the factor $Q_{\beta_j}^{-1/2}$, and differentiating the
exponential produces $t_j$ by the transverse triple-product identity.
Explicitly, on the absolute auxiliary torus in a fixed band,
\[
 \mathscr D_R=\partial_R+
 \sqrt Q\,d_r r^{d_r-1}\mathbf v_r\cdot\partial_Y,
 \qquad r=\sqrt Q R,\qquad \mathscr D_Z^*=\varepsilon\partial_Z.
\]
For the actual coefficient $c=(c_r,c_\theta,c_z)$, which has
$\partial_\theta c=0$, the operator is exactly
\[
 \operatorname{curl}_{*,\mathrm{coeff}}c
 =\bigl(-\varepsilon\partial_Zc_\theta,
       \varepsilon\partial_Zc_r-\mathscr D_Rc_z,
       (\mathscr D_R+R^{-1})c_\theta\bigr).
\]
Coefficients originally on a band auxiliary torus are pulled back to
the absolute torus before applying this formula; this retains the
integer-covering chain derivative in $\partial_Yc$.
In the same fixed chart the exact axial derivative is
\begin{align}
 \partial_zv
 &=\sum_j\operatorname{Re}\left[
 Q_{\beta_j}^{-A-D}
 \bigl(\partial_{Z_j}b_j+
   ik_{\beta_j}m_j(\partial_{Z_j}\Phi_j)b_j\bigr)
 e^{ik_{\beta_j}m_j\Phi_j}\right].
 \label{nsmom:actual_derivative}
\end{align}
The band power is fixed during chart differentiation and
$\partial_z=Q_{\beta_j}^{-D}\partial_{Z_j}$; auxiliary phase evaluation
adds no axial derivative because of \eqref{nsmom:phase}. Thus
\eqref{nsmom:actual_amplitude}--\eqref{nsmom:actual_derivative} substitute
the actual retained pair and all its curl terms into the bounds.
For example, writing
$\|b\|_{e,\beta}^2=\int R^e\langle|b|^2\rangle\,dR$ in that band,
the triangle inequality and $|\operatorname{Re}z|\le|z|$ give
\begin{align}
 \|v\|_e
 &\le\sum_j Q_{\beta_j}^{(e+1)/4-A}\|b_j\|_{e,\beta_j},
 \label{nsmom:scale_bound}\\
 \|\partial_zv\|_e
 &\le\sum_j Q_{\beta_j}^{(e+1)/4-A-D}
 \|\partial_{Z_j}b_j+
 ik_{\beta_j}m_j(\partial_{Z_j}\Phi_j)b_j\|_{e,\beta_j}.
 \label{nsmom:scale_derivative_bound}
\end{align}
The exponent $(e+1)/4$ is exactly the square root of the radial measure
factor $r^e\,dr=Q^{(e+1)/2}R^e\,dR$. Distinct bands and harmonics
remain distinct in these sums. No limiting estimate as $q\downarrow0$
has been inferred from the finite compact-patch norms.

For a single fixed chart, put
$F_e=Q^{-2A-1/2}E_e^*$ as in the covariance bridge and
$S^*=Q^{2A}S$. Direct change of radial variable gives
\[
 M_e=Q^{e/2-2A}M_e^*,\qquad
 M_e^*=\int R^eE_e^*\,dR,
 \qquad
 M_e^*=\varepsilon\partial_Z\int R^e S^*_{z,\tau_e}\,dR,
 \quad\varepsilon=Q^h,
\]
where $\tau_2=\theta$ and $\tau_1=z$. Indeed the axial flux has
factor $Q^{(e+1)/2-2A}$, its axial derivative contributes $Q^{-D}$,
and dividing by $Q^{e/2-2A}$ leaves $Q^{1/2-D}=Q^h$.
This proves the precise relation to the source chart moment without
altering its axial anisotropy.

\subsection{Exact relation to moments after phase evaluation}
\label{nsmom:evaluated}
Let
\[
 \mathcal S=\langle w_*\otimes v+v\otimes w_*+v\otimes v\rangle_\theta,
 \qquad \mathcal S^\circ=\mathcal S-\langle\mathcal S\rangle_Y.
\]
Thus $\langle\mathcal S\rangle_Y=S$ and
$\langle\mathcal S^\circ\rangle_Y=0$. For the evaluated physical
fields, denote the angular-only residual moments by $M_e^{\rm ev}$.
The proof of \eqref{nsmom:mean_zero} holds after evaluation. The
physical divergence, vector Laplacian and gradient identities therefore
prove, without any auxiliary averaging,
\[
 M_e^{\rm ev}
   =\partial_z\int_0^\infty r^e
        \mathcal S_{z,\tau_e}(r,z,t,Y(r,t))\,dr.
\]
The physical radial boundary terms still vanish by compact support;
the chain rule is included in the physical radial derivative before
that integration. Subtracting \eqref{nsmom:M2}--\eqref{nsmom:M1} proves
the exact comparison
\begin{equation}
 M_e^{\rm ev}-M_e
 =\partial_z\int_0^\infty r^e
     \mathcal S^\circ_{z,\tau_e}(r,z,t,Y(r,t))\,dr.
 \label{nsmom:evaluation_difference}
\end{equation}
For example its absolute value is bounded by
$\int r^e\sup_Y|\partial_z\mathcal S^\circ_{z,\tau_e}|\,dr$,
because $Y(r,t)$ is independent of $z$. Zero auxiliary Haar average
does not make the right side of \eqref{nsmom:evaluation_difference}
zero at the physical phase. The displayed map supplies precisely the
remaining auxiliary-oscillation flux; it is not silently identified
with the source fast mean.
\clearpage\part{Quantitative source-chart orders for the retained third seed}
\section{Finite quantitative source-chart orders for the retained third seed}
\label{qext:setup}
Fix one retained band, with
\[
0<Q\le1,\quad\varepsilon=Q^h,\quad A=\frac12+h,\quad D=\frac12-h,
\qquad r=\sqrt Q\,R,\quad z=Q^D Z .
\]
Here \(A_0>0\), \(\lambda _0\), and \(\mu=\lambda _0\) retain their original
values.  The source clock \(\nu=\sqrt{A_0}\) is distinct from the physical
viscosity \(\nu_{\rm NS}>0\).  The actual third seed and target are
\[
S_3=\frac{A_0}{\mu}\widehat\lambda _3^{-k_3-6},\quad
\Theta_3(t_3^{\rm in})=-S_3,\quad
-\Theta_3(t_3^a)=S_3e^{L_3}
=\frac{A_0}{\mu}\widehat\lambda _3^{-7/8},\quad
L_3=(k_3+\tfrac{41}{8})\log\widehat\lambda _3 .
\]
The last equality follows by adding the exact exponents
\(-k_3-6+k_3+41/8=-7/8\).  All constants below are finite maxima along
the actual \((\Theta_3,\Omega_3,\zeta_3)\) coefficient path and its actual
parent matrices; no endpoint or sign substitution is made.

\subsection{Physical curl and axial derivative}
For each retained harmonic let
\[
t_j=\chi_jB_jL_jy_3,\quad
c_j=\frac{i\,n_{\Phi,j}\times t_j}{k_{\beta_j}m_j|n_{\Phi,j}|^2},\quad
b_j=t_j+\operatorname{curl}_{*,\mathrm{coeff}}c_j,
\]
where the coefficient curl includes every cutoff, auxiliary chain,
cylindrical \(R^{-1}\), and evolving-path derivative.  Put
\[
v_j=\Re(Q^{-A}b_je^{ik_{\beta_j}m_j\Phi_j}),\quad
d_j=\partial_{Z_j}b_j+ik_{\beta_j}m_j(\partial_{Z_j}\Phi_j)b_j .
\]
For \(e=1,2\), put \(\alpha_e=(e+1)/4-A\). With \(\|a\|_e^2=\int R^e\langle|a|^2\rangle\,dR\),
\[
\|v_j\|_e\le Q^{(e+1)/4-A}\|b_j\|_e,\qquad
\|\partial_zv_j\|_e\le Q^{(e+1)/4-A-D}\|d_j\|_e .
\]
The powers are exact chart powers; the real-part fields obey these finite
upper bounds (a single pure nonzero harmonic has the sharper factor
\(2^{-1/2}\)).  Thus
\[
\begin{array}{c|cc}
e&\text{factor multiplying }\|b_j\|_e&\text{factor multiplying }\|d_j\|_e\\ \hline
2&Q^{1/4-h}&Q^{-1/4}\\
1&Q^{-h}=\varepsilon^{-1}&Q^{-1/2}
\end{array}
\]
as the factors multiplying the actual coefficient norms in the preceding bounds, before summing the finite retained harmonics.

\subsection{Flux products and moments}
Write separate old-parent orders
\[
\|w_{*,\tau}\|_e\le Q^{\omega_{e,\tau}}W_{e,\tau},\qquad
\|\partial_zw_{*,\tau}\|_e\le Q^{\dot\omega_{e,\tau}}\dot W_{e,\tau},
\quad \tau\in\{z,\theta\}.
\]
Put \(B_{j,e,\tau}=\|b_{j,\tau}\|_e\) and
\(D_{j,e,\tau}=\|d_{j,\tau}\|_e\).  Cauchy--Schwarz in the original
\(r^e dr\,d\theta\,dY\) measure gives the separate old--new terms
\[
\begin{aligned}
|J_{2,\mathrm{old}\times v}|&\le Q^{\alpha_2}\bigl(Q^{\omega_{2,z}}W_{2,z}B_{j,2,\theta}
 +Q^{\omega_{2,\theta}}W_{2,\theta}B_{j,2,z}\bigr),\\
|M_{2,\mathrm{old}\times v}|&\le Q^{\dot\omega_{2,z}+\alpha_2}\dot W_{2,z}B_{j,2,\theta}
 +Q^{\omega_{2,z}+\alpha_2-D}W_{2,z}D_{j,2,\theta}\\
&\quad +Q^{\dot\omega_{2,\theta}+\alpha_2}\dot W_{2,\theta}B_{j,2,z}
 +Q^{\omega_{2,\theta}+\alpha_2-D}W_{2,\theta}D_{j,2,z},\\
|J_{1,\mathrm{old}\times v}|&\le 2Q^{\alpha_1+\omega_{1,z}}W_{1,z}B_{j,1,z},\\
|M_{1,\mathrm{old}\times v}|&\le 2Q^{\dot\omega_{1,z}+\alpha_1}\dot W_{1,z}B_{j,1,z}
 +2Q^{\omega_{1,z}+\alpha_1-D}W_{1,z}D_{j,1,z}.
\end{aligned}
\]
The factors (1,2) are finite majorants from
\(|\operatorname{Re}z|\le|z|\); for a single pure harmonic they may be
replaced by the exact Gram factors.  The new quadratic terms are
\[
|J_{2,\mathrm{quad}}|\le Q^{2\alpha_2}B_{j,2,z}B_{j,2,\theta},\quad
|M_{2,\mathrm{quad}}|\le Q^{2\alpha_2-D}(D_{j,2,z}B_{j,2,\theta}+B_{j,2,z}D_{j,2,\theta}),
\]
\[
|J_{1,\mathrm{quad}}|\le Q^{2\alpha_1}B_{j,1,z}^2,\qquad
|M_{1,\mathrm{quad}}|\le 2Q^{2\alpha_1-D}B_{j,1,z}D_{j,1,z}.
\]
Substituting \(A=\tfrac12+h\), \(D=\tfrac12-h\), the quadratic powers are
\[
J_2:Q^{1/2-2h},\quad M_2:Q^{-h}=\varepsilon^{-1},\qquad
J_1:Q^{-2h}=\varepsilon^{-2},\quad M_1:Q^{-1/2-h}.
\]
These are finite norm majorants, not asymptotic source theorems; old-parent
orders remain the explicit \(\omega\) and \(\dot\omega\) symbols above.
The old field and old derivative orders stay independent in all four cross terms. No equal-order cross-term table is inferred from the quadratic powers.

The evaluated-phase correction remains the exact map
\[
M_e^{\rm ev}-M_e=\partial_z\int r^e
 \mathcal S^\circ_{z,\tau_e}(r,z,t,Y(r,t))\,dr,\qquad
 \langle\mathcal S^\circ\rangle_Y=0 .
\]
Define the finite normalized derivative majorant
\[
C_e^\circ(Z,T)=\int_0^\infty R^e\sup_Y
 \left|\partial_Z\bigl[Q^{2A}(\mathcal S^\circ_{z,\tau_e})^Q\bigr](R,Z,T,Y)\right|\,dR .
\]
The direct physical change of variables gives the exact finite envelope
\[
|M_e^{\rm ev}-M_e|\le Q^{(e+1)/2-2A-D}C_e^\circ(Z,T)
 =Q^{e/2-2A+h}C_e^\circ(Z,T).
\]
The zero auxiliary mean therefore does not cancel the evaluated flux.  This
uses the already differentiated normalized chart stress and introduces no
unsupported product of two stress norms.

\subsection{Retained viscosity}
The viscous term is the physical \(\nu_{\rm NS}\Delta v_j\).  In this chart
its radial, axial and angular costs carry \(Q^{-1}\), \(Q^{-2D}=Q^{-1+2h}\),
and \(Q^{-1}R^{-2}\), respectively, including the cylindrical
\(2R^{-2}\partial_\theta v_r\) coupling.  Thus no source \(\varepsilon\)
or clock \(\nu=\sqrt{A_0}\) is substituted for \(\nu_{\rm NS}\).  These
finite identities propagate the actual signed third seed and coefficient
path through \(v,\partial_zv,J_2,J_1,M_2,M_1\), while making no assertion
about a uniform \(Q\downarrow0\) endpoint.
\label{qext:conclusion}

\end{document}
